\section{相关工作}

从“非线性被引入到模型的哪个层次”这一视角看，现有非线性张量分解方法大致可归纳为四类。第一类方法在保留低秩因子骨架的前提下，直接修改潜在线性预测到观测条目的映射方式，即将非线性放在条目生成层。较早的一类工作通过广义损失或广义似然扩展经典分解模型，例如 generalized CP decomposition \cite{hong2020generalized} 在保留低秩参数化的同时，以更一般的观测损失替代平方误差，使模型能够适应二值、计数和鲁棒建模等场景；Bayesian Poisson Tucker decomposition \cite{schein2016bayesian} 则在 Tucker 框架下引入 Poisson 观测模型与层次贝叶斯结构，用于计数型多维事件数据建模。进一步地，InfTucker \cite{xu2011infinite}、Distributed Flexible Nonlinear Tensor Factorization \cite{zhe2016distributed} 与 Streaming Nonlinear Bayesian Tensor Decomposition \cite{pan2020streaming} 通过核方法、高斯过程或其稀疏近似，在条目层引入更强的非参数非线性表达。总体而言，这一路线与本文最为接近，因为它们都试图在不完全放弃低秩张量骨架的情况下增强观测生成机制；但相应代价通常是更重的推断过程、更多的核或变分超参数，以及相对较弱的结构透明性。

第二类方法将非线性放在因子本身的生成、约束或组织方式中，其核心不再是改变“条目如何由因子产生”，而是改变“因子如何被构造”。这类方法通常通过深度先验、时序网络或几何流形假设，把传统线性因子空间弯曲为受约束的非线性表示空间。典型地，DeepTensor \cite{saragadam2024deeptensor} 令低秩因子由深度网络生成，利用深度先验的隐式正则化来弥补经典低秩模型在复杂信号结构上的表达不足；Neural Tensor Factorization \cite{wu2019neural} 以及 DeepCP \cite{liu2017deepcp} 则将嵌入、时间依赖和神经网络生成机制结合起来，以适应时序交互和带噪观测场景；Neural Additive Tensor Decomposition \cite{ahn2024neural} 试图在表达能力与可解释性之间取得折中，通过对潜在成分施加神经网络变换来增强稀疏张量建模能力。另一些工作则更强调几何结构，例如 3D head pose estimation 中的张量分解与非线性流形建模 \cite{derkach2019tensor}，通过对姿态因子施加显式流形约束来提升泛化与可解释性。相比条目层非线性，这类方法往往具有更高的建模自由度，但同时也更依赖网络结构设计或先验选择，理论分析和优化稳定性也更复杂。

第三类方法进一步将张量分解算子本身“网络化”，即把非线性直接嵌入到模乘、核收缩或层级分解结构内部，形成深度张量模型。其代表思想不是在传统低秩分解外部附加一个非线性部件，而是把 Tucker、TT 等结构重新组织为多层、可堆叠的非线性变换单元。比如，Stacked Tucker Decomposition \cite{xu2024stacked} 在模乘之后插入非线性激活，并通过层级堆叠构造深度化的 Tucker 结构；NMTucker \cite{varolgunes2023nmtucker} 在 Matryoshka 式 Tucker 框架中显式引入非线性组件，以改善高稀疏补全任务中的表达能力；Nonlinear Tensor Train \cite{wang2021nonlinear} 则在 TT 核收缩之间加入激活，用于深度网络压缩与表示学习。与前两类方法相比，这一路线的优势在于能够在保持张量参数共享结构的同时获得接近深度网络的表达能力，但其代价是模型越来越偏向“神经化分解器”，从而削弱了经典低秩模型在参数含义、结构解析和理论刻画上的直接性。

第四类方法并不一定直接修改张量分解本身，而是将张量视为表示外部非线性结构的紧凑载体，再通过分解实现高效建模与计算。这类方法常见于多项式系统、Volterra 系统、非线性变换域表示等场景，其中非线性结构通常先由物理机制、解析展开或前端变换给定，张量主要负责高阶系数的压缩表示与数值加速。典型如 TNMOR 和 STORM \cite{liu2015model,deng2016storm}，它们先对非线性电路或动力系统作多项式/Taylor/Volterra 展开，再将各阶系数重排为高阶张量并做 CP 或对称分解，从而构造保留高阶非线性的降阶模型；另一类工作如多维非线性变换张量表示 \cite{shi2025multidimensional}，则通过卷积网络和激活函数构造数据自适应的非线性变换域，再在该域内施加张量低秩约束。总体来看，这些方法说明张量不仅可以作为“非线性分解器”，也可以作为“非线性结构的压缩表达工具”；但在多数情形下，非线性的来源位于张量模型之外，而非在经典 Tucker 骨架内部形成一个受控、显式且可直接解释的条目级扩展。

与上述工作相比，NTD-PL 的定位更接近“保留 Tucker 骨架的条目生成机制扩展”，但又与广义似然、核方法、高斯过程和深度解码器等路线有所区别。本文并不依赖无限维核映射、重型概率推断或深度生成网络，而是在共享 Tucker 潜在张量之上引入显式、有限阶、逐元素的多项式链接，使非线性能力以少量系数受控进入模型。这一设计使得 NTD-PL 一方面继承了 Tucker 的低秩结构、参数效率与多模可解释性，另一方面又能够通过明确的阶数展开刻画条目层残差、给出与 Tucker 的嵌套关系以及相应的逼近分析。就现有文献而言，NTD-PL 的创新点正在于：它不是将非线性完全外包给核函数、神经网络或外部多项式系统，而是在经典 Tucker 框架内部构造了一个显式、可分析、可逐阶控制的非线性链接机制，从而在表达能力、结构透明性与优化可实现性之间取得新的折中。
